% r07 Corrections to King's papers used in the rotation part
\section{Corrections to King's papers used in the rotation part}

We list the points in~\cite{K1,K2} that affect the use of King's theorems in the rotation part. None of them affects the existence theorem \EC\ (\Cref{hyp:s03:EC}), the existence of King's cube limits at zero scalar source, which is proved in \cref{sec:k01} (\Cref{rot:rem:frame}). The point concerning the Wick square and the constant $\ln5$ was treated in \Cref{rot:rem:log}.

\subsection{The large-field term in Theorem~3.1 of~\cite{K1}}

Theorem~3.1 of~\cite{K1} compares the lattice generating functional on a torus $T$ with the small-field integral
$I^{(k)}\defeq\int(dA_k)(d\phi_k)\,\chi_k\exp[-S^{(k),1}]$ of the effective action after $k$ block-spin steps, $\eta\defeq5^k\eps_K$. As printed on p.~655 the upper bound reads
\[
Z^{\eps_K}(T_{\eps_K},g,h)\le\int(dA_k)(d\phi_k)\,\chi_k\exp\big[-S^{(k),1}+C\eta^\sigma|T|\big]+\exp\big[-p(\eta)^2+C|T|\big],\qquad(3.4)
\]
with $p(\eta)=b_0(1+\log\eta^{-1})^p$, followed by ``Once again, $C$ depends on $g$, $h$ and $|T|$''; the lower bound (3.3) has the small-field integral with the exponent $\mp C\eta^\sigma|T|$ (in a lower bound the sign must be $-$). The second term of (3.4) is an additive large-field remainder; it is used in~\cite{K1} only at fixed volume ((3.10)--(3.13) there), and in the same way in King's torus comparison~\cite[(4.7)]{K2}, where it is harmless.

\begin{proposition}[the large-field term]\label{rot:prop:largefield}
Suppose (3.3)--(3.4) hold with $C$ independent of $|T|$, and $I^{(k)}\ge e^{-C''|T|}$. Then
\[
-C\eta^\sigma\le|T|^{-1}\log\big(Z/I^{(k)}\big)\le\max\{C\eta^\sigma,\ C+C''-p(\eta)^2/|T|\}+|T|^{-1}\log2 .
\]
At fixed $\eta$ the right side tends to $C+C''$ as $|T|\to\infty$, not to $O(\eta^\sigma)$.
\end{proposition}

\begin{proof}
$Z/I^{(k)}\le e^{C\eta^\sigma|T|}+e^{-p(\eta)^2+(C+C'')|T|}$ by (3.4) and the lower bound on $I^{(k)}$, and $\log(x+y)\le\max(\log x,\log y)+\log2$ for $x,y>0$. The lower bound is (3.3).
\end{proof}

Hence the printed bounds do not give a two-sided representation of the free energy per unit volume by a small-field effective integral at a fixed coarse scale, uniformly in the volume, and arguments that need such a representation cannot be based on them. In the rotation part the volume is fixed wherever bounds of this type enter (the one-sided bounds \hypF\ of \Cref{rot:hyp:F} are posed on one torus at a time), and the large-field term is absorbed on the left by a polynomial lower bound for the walled partition function (\Cref{rot:rem:closing}). The constant of~\cite[Thm.~2.1(ii)]{K1} likewise depends on $|T|$, so King's theorems give no bound on torus functionals uniform in the torus; where such a bound is used (\Cref{rot:lem:K}(d)) it is taken on a fixed torus and comes from the existence of the torus limit.

\subsection{King's torus conjecture: Theorem~2.4 of~\cite{K2} as printed}

In~\cite{K1,K2} the curvature source is $F(g)=\sum_p\eps^3g(p)F(p)$ with one function $g$ on all plaquettes~\cite[(2.12)]{K2}, and in Theorem~2.4 of~\cite{K2} the source is rotated as a scalar. For a lattice $\Lambda$ with frame $(f_\alpha)$ descending to $\DJ$ (\Cref{rot:def:frames,rot:def:tori}) and a real function $g$ on $\DJ$ put $F^{sc}_\Lambda(g)\defeq\sum_p\eps^3g(c(p))F(p)$, the sum over the plaquettes $p=(x;\alpha,\beta)$ of one fundamental cell, $c(p)$ their centres.

\begin{proposition}[the scalar reading]\label{rot:prop:scalar}
Let $e=0$, $R\defeq R_{12}(\thz)$, $\ell\defeq50MJ$ and $g(x)\defeq\cos(2\pi x_1/\ell)$. Then
\[
\lim_{k\to\infty}\big\langle e^{F^{sc}_{\Lk}(g)}\big\rangle_{\Lk,\DJ}\ne\lim_{k\to\infty}\big\langle e^{F^{sc}_{R\Lk}(g)}\big\rangle_{R\Lk,\DJ},
\]
where $R\Lk$ carries the frame $(Re_\alpha)$ and descends to $\DJ$ by \Cref{rot:lem:geom}(b).
\end{proposition}

\begin{proof}
At $e=0$ the photon is the centred Gaussian with covariance $\delta_{\alpha\beta}(-\Delta_\eps+\lmu^2)^{-1}$ on the torus, so $\langle e^{F^{sc}(g)}\rangle=e^{\frac12\Var F^{sc}(g)}$. On a lattice with frame $(f_\alpha)$, $\sum_{\alpha<\beta}F(f_\alpha,f_\beta)=\ip uB$ for a continuum field strength $F=dA$, $B=\star F$, where $u\defeq\sum_{\alpha<\beta}f_\alpha\times f_\beta$. For the plane wave $g$ only the momenta $\pm k_0$, $k_0=(2\pi/\ell)e_1$, contribute, and the lattice symbols converge, so
\[
\lim_{k\to\infty}\Var F^{sc}(g)=\frac{\ell^3}2\,\frac{|k_0|^2\,(|u|^2-u_1^2)}{|k_0|^2+\lmu^2}.
\]
For $\Lk$, $u=(1,-1,1)$ and $|u|^2-u_1^2=2$; for $R\Lk$, $u=R(1,-1,1)=(\frac75,-\frac15,1)$ and $|u|^2-u_1^2=\frac{26}{25}$. The limiting variances differ by the factor $\frac{25}{13}$.
\end{proof}

With $2$-form sources sampled in each lattice's own frame (\Cref{rot:def:frames}) the limit on either lattice is the variance of $\int\ip GF$, which does not depend on the frame, and the obstruction disappears. This reading is consistent with King's proof, which treats the source components separately (\Cref{rot:rem:frame}); it is the form of the torus comparison \Kz, \Ka\ (\Cref{rot:hyp:K,rot:hyp:Ka}). The second error in the printed statement concerns the Wick constant (\Cref{rot:rem:log}). Since the printed statement is false on both counts, we refer to it as King's torus conjecture; what the rotation part uses is the torus comparison \Kz\ (\Cref{rot:hyp:K}), in the corrected form. We note also that Theorem~2.4 is not, by itself, a statement about King's infinite-volume state: the passage from tori to the state~\cite[(2.24)--(2.25)]{K2} is explicitly conjectural there (``we can expect that a cluster expansion would give the following result''), and the rotation part does not use it (\Cref{rot:rem:shape}).

\subsection{The bound (4.33) of~\cite{K1}}

In the proof of~\cite[Lemma~4.5]{K1} (p.~674) the positivity (4.33) is obtained through (4.36)--(4.37), which treat the block average as a Fourier multiplier; that step is false as printed. The conclusion holds. We use King's setting at the unit scale of step $k$; in this subsection $a$ and $a_j$ are King's coefficients of~\cite[(4.5)]{K1}, not a scheme, and $\eta$ is King's unit scale, not a tolerance. $\Omega$ is a periodic box of the unit lattice $\Z^3$ that is a union of blocks of side $5$, the background vanishes, $\Delta^{(k)}$ and $\Delta^{(k+n)}$ are the translation-invariant operators of the flow on $\Omega$, $Q$ is the block average, so that $Q^*Q=P$ is the orthogonal projection onto block-constant functions, and $-\Delta$ is the lattice Laplacian of $\Omega$. King's coefficients satisfy $a_j\ge\frac{24}{25}a$, and $\Delta^{(j)}(p)^{-1}=a_j^{-1}+\sum_l|u(p+l)|^2/\Delta^{\eta}(p+l)$~\cite[(4.5)]{K1}, where $u$ is the symbol of the average over a unit block of the fine lattice of spacing $\eta$, $\sum_l|u(p+l)|^2=1$, the sum runs over $l\in2\pi\Z^3$ with $p+l$ in the Brillouin zone of the fine lattice, and $\Delta^\eta(q)=\sum_\mu4\eta^{-2}\sin^2(\eta q_\mu/2)+m_k^2$ with $m_k^2\ge0$.

\begin{proposition}[the bound (4.33)]\label{rot:prop:433}
For $s\in[0,1]$, all $k,n$ and $m_k^2\ge0$, the operator $X(s)\defeq s\Delta^{(k)}+(1-s)\Delta^{(k+n)}+\frac a{25}Q^*Q$ on $\ell^2(\Omega)$ satisfies $X(s)\ge\gamma_0$, where
\[
\gamma_0\defeq\min\Big(4c_\Delta\sin^2\frac\pi{10},\ \frac a{25}\Big),\qquad
c_\Delta\defeq\Big(\frac{25\cdot3\pi^2}{24a}+\frac{\pi^2}4\Big)^{-1}.
\]
\end{proposition}

\begin{proof}
(i) For $|t|\le\pi$, $4\sin^2(t/2)\ge4t^2/\pi^2$, so $\Delta^\eta(q)\ge4|q|^2/\pi^2$ on the fine Brillouin zone; and $|p_\mu+l_\mu|\ge|p_\mu|$ for $|p_\mu|\le\pi$, $l_\mu\in2\pi\Z$. With $\sum_l|u(p+l)|^2=1$ this gives $\Delta^{(j)}(p)^{-1}\le a_j^{-1}+\frac{\pi^2}4|p|^{-2}$, and since $|p|^2\le3\pi^2$ and $a_j\ge\frac{24}{25}a$, $\Delta^{(j)}(p)\ge c_\Delta|p|^2\ge c_\Delta\sum_\mu4\sin^2(p_\mu/2)$, the symbol of $-\Delta$. So $s\Delta^{(k)}+(1-s)\Delta^{(k+n)}\ge c_\Delta(-\Delta)$.
(ii) On one block $B=\{0,\dots,4\}^3$ the Neumann Laplacian is a sum of three path-graph Laplacians, whose smallest non-zero eigenvalue is $\lambda_1=4\sin^2\frac\pi{10}$; hence $\sum_{x\in B}|\phi(x)-\phi_B|^2\le\lambda_1^{-1}\sum_{\{x,x'\}\subset B}|\phi(x)-\phi(x')|^2$, the sum over nearest-neighbour pairs, $\phi_B$ the block mean. Summing over the blocks, whose inner bonds are bonds of $\Omega$, gives $\|(1-P)\phi\|^2\le\lambda_1^{-1}\langle\phi,-\Delta\phi\rangle$.
(iii) Therefore $\langle\phi,X(s)\phi\rangle\ge c_\Delta\lambda_1\|(1-P)\phi\|^2+\frac a{25}\|P\phi\|^2\ge\gamma_0\|\phi\|^2$. For the photon the same holds componentwise.
\end{proof}

\subsection{Minor points}

\begin{remark}[minor points]\label{rot:rem:minor}
(i) In~\cite[(3.41)]{K2} the real part of a scalar source is split as $h_+-h_-$ with $h_\pm\ge0$ and no further regularity stated; the argument of \Cref{rot:thm:full}(iii) uses smooth non-negative parts.
(ii) The inequalities~\cite[(3.24)--(3.29)]{K2} are stated for ratios of expectations of $e^{r^2(h)+\mathrm iA(g)}$; the rotation part uses them in the form of~\cite[Thms.~6.1, 6.2]{BFS1}, \Cref{lem:s02:BFS}, which also covers killed sites and the two-level walls.
(iii) King's bound~\cite[(3.43)]{K2} is asserted for the Neumann--Dirichlet cubes by ``a standard argument (see [GJ1])'', with no argument printed; it is not a bound on torus functionals, which lie above the cube functionals on the cone (\Cref{rot:lem:cutwall}). The torus moment constant of \Cref{rot:lem:K}(d) is therefore taken at a fixed torus and comes from the existence of the torus limit under \Ka, not from (3.43).
\end{remark}
